When @@M0@@ is of characteristic zero, i.e. its structure sheaf is a sheaf of @@M1@@-algebras, the problem is solved simply by taking
\sgaStableTarget{sga6.e14.eq.1-2}@@M2@@
where the @@M3@@-th tensor power (in the sense of derived categories) of @@M4@@ is regarded as an object of @@M5@@, with @@M6@@ the group algebra of the symmetric group @@M7@@ with coefficients in @@M8@@, and where the superscript alt denotes passage to the ``alternating part under the action of @@M9@@'' (an operation that makes sense by the characteristic-zero hypothesis). A formula of this type was already used in [8] (cited below as [RRR]) for the first proof of the Riemann--Roch theorem.

Without a hypothesis on the characteristic, when @@M10@@ is zero in degrees @@M11@@, Dold--Puppe [6] shows that @@M12@@ is canonically homotopic to the complex associated with a semi-simplicial Module @@M13@@ canonically associated with @@M14@@, and a fairly natural candidate for @@M15@@ (for @@M16@@ flat, say) would be the complex associated with the semi-simplicial Module @@M17@@ (whose component of degree @@M18@@ is @@M19@@). But to extend the definition to complexes of @@M20@@ no longer assumed to be zero in degrees @@M21@@, one would have to study the behaviour of the preceding operation with respect to the translation operation @@M22@@ (or \uline{suspension}, in the terminology of [6]); unfortunately, there does not seem to be an isomorphism @@M23@@ that would make the desired extension trivial.

\sgaStableTarget{sga6.e14.sec.2}\section*{2. The Riemann--Roch formula without projectivity hypotheses}

\sgaStableTarget{sga6.e14.para.2-1}\textbf{2.1.} Once a natural @@M24@@-structure on the @@M25@@ is available, it is possible to formulate the Riemann--Roch theorem, in the form considered in Exp. VIII, for any proper complete-intersection morphism @@M26@@ (and even, more generally, for a relative scheme @@M27@@ over a locally ringed topos @@M28@@, with @@M29@@ still subject to the corresponding relative conditions). The Todd class of @@M30@@ is defined by means of the class in @@M31@@ of the relative cotangent complex @@M32@@ for @@M33@@, constructed in Exp. 0 4.4; this complex is perfect here by the local hypothesis on @@M34@@. It seems clear that a proof of the Riemann--Roch formula in this general case will require the introduction of essentially new ideas, and it is possible that these same ideas could provide the key
